\begin{frame}{Conclusions and questions}
  \framesubtitle{Q1--Q7 | SUMMARY}
  \vspace{0.08cm}
  \begin{columns}[T,onlytextwidth]
    \begin{column}{0.65\textwidth}
      \begin{enumerate}
        \item \textbf{Q1--Q2:} sublinear regret means vanishing regret per step; credible finite-horizon evidence needs several complementary diagnostics.
        \item \textbf{Q3:} polynomially decaying exploration gives the clearest sublinear trend; a low average can conceal rare lock-in.
        \item \textbf{Q4--Q5:} UCB \(c=2\) estimates arm values best, while UCB \(c=1\) achieves the lowest regret.
        \item \textbf{Q6--Q7:} gradient bandit works well only in a moderate step-size range; aggressive updates can create bimodal outcomes.
      \end{enumerate}
      \vspace{0.20cm}
      \smallnote{All reported comparisons use the same fixed ten-arm Gaussian testbed and notebook protocols. Claims are empirical unless stated otherwise.}
    \end{column}
    \begin{column}{0.31\textwidth}
      \centering
      \vspace{0.42cm}
      {\color{Teal}\rule{0.82cm}{2.6pt}}\par
      \vspace{0.24cm}
      {\fontsize{25}{28}\selectfont\bfseries\color{Navy}Questions?}\par
      \vspace{0.18cm}
      {\fontsize{10.2}{12.4}\selectfont\color{Muted}5-minute discussion}
    \end{column}
  \end{columns}
\end{frame}
